\chapter{实验结果与分析}
\echapter{Experimental Results and Analysis}

\section{实验设置与评价指标}
\esection{Experimental Setup and Metrics}

本文实验主要围绕两个环境展开。其一是原型协作环境 \texttt{join1}，用于快速验证通信机制是否具有基础有效性；其二是 SMAC 基准中的 \texttt{5m\_vs\_6m} 任务，用于评估算法在部分可观测协作微操场景中的实际表现。

为便于对不同方法进行统一比较，本文主要采用以下三个指标：
\begin{enumerate}
    \item 峰值胜率（peak），即测试胜率曲线中的最大值，用于衡量方法上限；
    \item 最终胜率（final），即最后一个有效测试点的胜率，用于衡量训练结束时的结果；
    \item 末 10 点平均胜率（last10），即最后 10 个测试点的平均值，用于衡量方法后期稳定性。
\end{enumerate}

此外，在机制诊断阶段，本文引入以下中间指标用于分析通信模块的真实行为：
\begin{enumerate}
    \item 攻击注意力熵（attack attention entropy），衡量攻击通信路由的选择性，接近 $\ln(k)$ 表示对 $k$ 个候选的均匀关注；
    \item 攻击无通信概率（attack no-comm probability），衡量攻击流选择沉默的频率；
    \item 攻击门控均值（attack gate mean），衡量攻击通信在决策中的实际融合权重；
    \item 通信隔离评估（communication isolation eval），通过对同一 checkpoint 分别执行正常评估、强制开门、强制关门和禁用攻击通信四种模式，直接测量攻击通信的边际贡献。
\end{enumerate}

\section{\texttt{join1} 环境中的通信原型验证}
\esection{Communication Verification in \texttt{join1}}

在正式进入 SMAC 任务之前，本文首先在 \texttt{join1} 环境中检验通信机制的基础有效性。图 \ref{fig:join1} 展示了 \texttt{qmix}、受预算约束的稀疏 MAPPO 方法，以及两组 MAIC 类通信方法在该环境中的测试胜率曲线。

\begin{figure}[htbp]
    \centering
    \includegraphics[width=0.88\textwidth]{generated/join1_test_winrate_comparison.png}
    \caption{\texttt{join1} 环境下不同方法的测试胜率对比}
    \label{fig:join1}
\end{figure}

从图中可以看出，受约束的稀疏通信在该环境中具有明显优势，而缺乏有效结构约束的通信方法则可能出现几乎完全失效的情况。这一现象说明，通信机制并不是天然有效的协作增强器；消息编码、路由与融合方式若设计不当，反而会直接破坏训练过程。

\section{SMAC \texttt{5m\_vs\_6m} 上的 Backbone 构建与训练动力学}
\esection{Backbone Construction and Training Dynamics on SMAC \texttt{5m\_vs\_6m}}

\subsection{Backbone 基线构建}

在 \texttt{5m\_vs\_6m} 任务上，本文首先围绕 vanilla MAPPO 主线构建可靠主干基线（Backbone）。经过集中式 Critic（centralized critic）、active masks、proper time-limit、学习率衰减与 KL 约束等实现细节的对齐与修正，vanilla MAPPO 能够在该任务上稳定学习。

与此同时，实验发现 Backbone 训练存在显著的种子敏感性。在早期配置（entropy coefficient $c_e=0.02$）下，3 个种子的最终胜率分别为 $0.575$、$0.025$ 和 $0.025$，仅 seed 1 成功。通过将 $c_e$ 提升至 $0.03$ 并降低 critic 学习率至 $5\times10^{-4}$，成功率提升至 2/3（seed 1: $0.325$、seed 2: $0.025$、seed 3: $0.200$，均为 $2.5\times10^{6}$ 步处的测试胜率）。

\subsection{崩溃-恢复动力学}

进一步分析 Backbone 的训练曲线后，本文发现一个重要的训练动力学特征：成功种子的学习曲线普遍经历“初期上升→中期崩溃→后期恢复并加速”三个阶段。图 \ref{fig:backbone_dynamics} 展示了基线形成与 warm-start 对照，并为后续通信 fine-tune 的起点选择提供依据。

\begin{figure}[htbp]
    \centering
    \includegraphics[width=0.95\textwidth]{generated/baseline_officialish_vs_warmstart.png}
    \caption{vanilla MAPPO Backbone 的形成过程与 warm-start 对照}
    \label{fig:backbone_dynamics}
\end{figure}

该现象的深层含义是：Backbone 本身的训练动力学非常复杂，存在多个训练可到达的策略吸引子，种子决定了其落入哪个吸引子。这对后续通信实验提出了关键约束——如果 Backbone 尚未完成崩溃-恢复循环并进入稳定平台期，就引入通信 fine-tune，那么通信的“增值”将无法与 Backbone 自身的持续改善区分开来。

基于这一发现，本文后续所有通信实验均采用 warm-start 策略：先在无通信条件下训练得到强 Backbone 检查点（约 $0.675$ 胜率），再从该检查点出发做通信 fine-tune。该策略同时规避了冷启动通信直接失败的问题（详见后续分析）。

\subsection{冷启动通信的结构性失败}

本节对冷启动（from-scratch）通信训练的系统性诊断实验进行总结。实验围绕 8 组递进式配置展开，涵盖信息增益门控（IGG）、固定门控、warm-up 门控和门控退火等多种策略，共计超过 70 次完整训练运行。核心发现如下。

\textbf{IGG 自动关门。} 在冷启动场景下，IGG 机制准确识别出通信对训练初期的策略是有害的，因此将攻击门控和移动门控均压制在约 $0.02$--$0.04$ 的地板值附近。虽然该版本获得了冷启动通信中的最高胜率（约 $11.25\%$），但门控从未打开——通信实质上被完全关闭了。此时 Backbone 仅以微量残留通信噪声干扰的形式进行局部策略学习。

\textbf{固定门控摧毁 Backbone。} 若强制将门控固定在 $g=0.5$ 并移除所有 IGG 复杂度（Plan G 实验），Backbone 完全无法学习，胜率仅为约 $1.25\%$。这证明通信架构本身（注意力路由 + delta 融合）才是冷启动失败的根本噪声源，而非门控学习机制。

\textbf{Warm-up 带来短暂峰值但不可持续。} 通过通信 warm-up 因子从 $0$ 线性增加到 $1$（B1 实验），Backbone 在保护期内建立了一定基础，峰值胜率达到 $15\%$——超过了 IGG 策略。但 warm-up 结束后，固定门控 $g=0.5$ 的强通信噪声在约 $4\times10^{5}$ 步后导致训练崩溃，胜率回落至 $0\%$。

\textbf{门控-学习率同构性。} 综合分析表明，门控值在功能上等价于“通信学习率”：
\[
\frac{\partial L}{\partial\theta_{\text{comm}}} \propto g \cdot \frac{\partial L}{\partial\tilde{\ell}} \cdot \frac{\partial\Delta}{\partial\theta_{\text{comm}}}.
\]
大门控（如 $g=0.5$）提供强 delta 梯度，加速协调学习但积累噪声，最终导致训练发散；小门控（如 $g\approx0.04$）梯度弱，学习缓慢但训练稳定。基于这一洞察，本文提出了门控退火（Gate Annealing）策略：将门控视为第二个学习率，早期使用大门控快速学习，后期退火至小门控稳定收敛。

\textbf{Hidden-State Fusion 验证。} 为从根本上解决“通信不改变决策”的问题，本文进一步尝试将通信注入点从 logit 层移至 hidden state 层：
\[
h_i^{\text{fused}} = (1-g) \cdot h_i^{\text{local}} + g \cdot h_i^{\text{comm}}, \quad
\tilde{\ell}_i = f_{\pi}(h_i^{\text{fused}}).
\]
实验中的 cosine 相似度从 $1.0$ 持续下降至 $0.816$，验证了 hidden fusion 确实让通信改变了策略表示。然而该机制并未转化为胜率提升（峰值仅 $3.75\%$，低于 C2 的 $12.5\%$），说明单层线性 policy head 难以从正交方向的信息中提取有用决策信号。

综合上述发现，冷启动通信在现有架构下始终无法产生稳定正向收益。这直接促成了本文后续的两阶段策略设计。

\section{五类轻量通信算法对比}
\esection{Comparison of Five Lightweight Communication Algorithms}

在固定强 Backbone 后，本文依次实现并比较了五类轻量通信算法。图 \ref{fig:lightweight_comm_summary} 展示了不同算法及其代表配置的峰值、最终点与末 10 点平均统计，图 \ref{fig:lightweight_comm_curves} 则给出了若干代表性 run 的测试胜率曲线。

\begin{figure}[htbp]
    \centering
    \includegraphics[width=0.96\textwidth]{generated/lightweight_comm_algorithm_summary.png}
    \caption{五类轻量通信算法各代表配置的统计结果对比}
    \label{fig:lightweight_comm_summary}
\end{figure}

\begin{figure}[htbp]
    \centering
    \includegraphics[width=0.92\textwidth]{generated/lightweight_comm_representative_curves.png}
    \caption{五类轻量通信算法代表性 run 的测试胜率曲线}
    \label{fig:lightweight_comm_curves}
\end{figure}

为便于总结，表 \ref{tab:main_results} 给出了各阶段代表配置的关键统计结果。

\begin{table}[htbp]
    \centering
    \caption{各阶段代表配置的关键结果对比}
    \label{tab:main_results}
    \begin{tabular}{|p{3.5cm}|p{4.0cm}|c|c|c|}
        \hline
        算法类别 & 代表实现 & peak & final & last10 \\
        \hline
        基线 Backbone & warm-start 微调对照组 & 0.854 & 0.625 & 0.642 \\
        \hline
        最小侵入通信 & 端到端残差接入 & 0.813 & 0.646 & 0.563 \\
        \hline
        锐化路由通信 & 软锐化路由 & 0.854 & 0.667 & 0.583 \\
        \hline
        动作意图共享 & 端到端动作意图共享 & 0.938 & 0.563 & 0.600 \\
        \hline
        攻击子空间融合 & 基础目标融合 & 0.813 & 0.750 & 0.613 \\
        \hline
        攻移双流通信 & Top-1 移动路由 + Softplus 门控 & 0.771 & 0.688 & 0.588 \\
        \hline
    \end{tabular}
\end{table}

综合图表与表格结果，可以得到以下观察。

首先，最小侵入通信证明了通信分支能够在强 Backbone 上稳定接入，而不会立即破坏原有策略。其次，锐化路由通信在注意力锐化后整体表现优于最小侵入通信，说明通信模块确实受益于更强的路由选择性。再次，动作意图共享通过共享动作意图，将通信内容从抽象隐层特征转为具有明确物理意义的消息，从而获得了当前最强的原始峰值表现。随后，攻击子空间融合将通信严格限制在攻击 logits 子空间中，使通信模块的作用更可解释，并取得较好的最终收尾表现。最后，攻移双流通信则表明攻击协同与移动协同应区别对待。

\section{攻击通信流的机制修复与隔离评估}
\esection{Attack Stream Repair and Isolation Evaluation}

\subsection{攻击注意力死亡问题}

在攻移双流通信的基础上，诊断实验揭示了一个严重问题：攻击通信流的注意力熵在所有 warm-start 配置中均为 $0.0$，表明攻击注意力完全饱和——始终硬选同一个路由目标。这与移动流的健康熵（约 $0.4$--$0.5$）形成鲜明对比。问题根因在于攻击流使用 $\text{top-}k=1$ 时，softmax 退化为退化的 one-hot 分布，且没有任何梯度信号能够告知注意力“应该选谁”。

为解决该问题，本文对攻击通信流进行了系统性的机制修复，所有实验均基于 warm-start 从同一 Backbone 检查点出发。实验结果汇总于表 \ref{tab:attack_repair}。

\begin{table}[htbp]
    \centering
    \caption{攻击通信流机制修复的关键指标对比（3 seed 平均）}
    \label{tab:attack_repair}
    \begin{tabular}{|l|c|c|c|c|}
        \hline
        配置 & 攻击熵 & 无通信概率 & 攻击门控 & 平均峰值胜率 \\
        \hline
        原始 v6（top1, 无约束） & 0.000 & 0.01--0.02 & 0.06--0.07 & 0.629 \\
        \hline
        top2 only（无正则） & 0.693 & 0.87 & 0.06 & -- \\
        \hline
        top2 + selective entropy & 0.350 & 0.87 & 0.06 & -- \\
        \hline
        top2 + silence budget & 0.429 & 0.011 & 0.061 & 0.704 \\
        \hline
        top2 + silence + align & 0.449 & 0.010 & 0.067 & 0.746 \\
        \hline
    \end{tabular}
\end{table}

修复过程分为以下步骤。

\textbf{步骤 1：top2 软路由。} 将攻击通信的 top-$k$ 从 $1$ 改为 $2$，使攻击流在 2 个队友之间进行软分配。该步骤成功将攻击注意力熵从 $0$ 提升至约 $0.693$（$\ln(2)$），但注意力退化为“均匀听两人”——仍然没有学到有意义的路由选择。

\textbf{步骤 2：upper-only 选择性熵正则。} 引入攻击注意力熵的 upper-only 正则化：仅当熵高于目标值 $0.50$ 时才施加惩罚，低于目标时不干预。该步骤在 top2 only 的退化结果（C3 实验中 attention entropy 贴 $0.693$、no-comm 概率接近 $0$）基础上，成功将熵压制至约 $0.35$。然而，攻击流发现了一条更简单的路径——通过选择“无通信”token 来满足正则约束。此时攻击无通信概率飙升至 $87.5\%$，攻击流几乎完全沉默。

\textbf{步骤 3：沉默预算。} 引入攻击无通信评分惩罚（silence budget, penalty $=0.3$），使选择 no-comm token 需要付出代价。该步骤成功将攻击无通信概率从 $87.5\%$ 压至约 $1\%$，攻击熵稳定在 $0.35$--$0.44$ 的理想区间，三 seed 平均峰值胜率从 v6 原版的 $0.629$ 提升至 $0.704$，且跨 seed 方差从 $20\%$ 缩小至 $8.75\%$。

\textbf{步骤 4：意图对齐。} 在 silence budget 基础上引入攻击意图对齐损失（alignment loss），鼓励攻击通信的消息内容与接收方的实际攻击目标一致。该步骤将三 seed 平均峰值胜率进一步提升至 $0.746$，且末 10 点平均胜率达到 $0.661$——远超前序所有版本。

\subsection{通信隔离评估}

为直接测量攻击通信在执行期的真实边际贡献，本文设计了通信隔离评估实验。对同一检查点分别执行 128 局测试，对比以下四种模式：
\begin{enumerate}
    \item 正常模式（normal）：攻击通信保持训练时的状态；
    \item 强制开门（gate\_open）：攻击门控强制设为 $1.0$；
    \item 强制关门（gate\_closed）：攻击门控强制设为 $0.0$；
    \item 禁用攻击通信（no\_attack）：完全关闭攻击通信流。
\end{enumerate}

对最优 align 配置（三 seed 平均峰值 $0.746$）的三个末期检查点进行评估，结果如表 \ref{tab:comm_isolation} 所示。

\begin{table}[htbp]
    \centering
    \caption{攻击通信隔离评估结果（align 配置）}
    \label{tab:comm_isolation}
    \begin{tabular}{|l|c|c|c|c|c|}
        \hline
        Seed & normal & gate\_open & gate\_closed & no\_attack & $\Delta$ (normal $-$ no\_attack) \\
        \hline
        Seed 1 & 0.383 & 0.422 & 0.383 & 0.383 & 0.000 \\
        Seed 2 & 0.750 & 0.766 & 0.750 & 0.750 & 0.000 \\
        Seed 3 & 0.711 & 0.656 & 0.664 & 0.664 & $+0.047$ \\
        \hline
    \end{tabular}
\end{table}

评估结果揭示了以下几个层次的重要发现。

\textbf{通信机制的直接证明。} Seed 3 上 $+4.7\%$ 的正向边际贡献（normal $=$ $71.1\%$ vs no\_attack $=$ $66.4\%$）是攻击通信有效性的最直接证据。该增益并非来自随机波动——四组独立评估（各 128 局）之间的系统差异清楚表明，攻击通信在该检查点上确实为集火协同提供了执行期价值。

\textbf{Backbone 质量的约束作用。} Seed 1 的评估胜率（$38.3\%$）远低于训练期间的测试胜率（约 $68\%$），且四种模式结果一致。通信 fine-tune 过程在该种子上破坏了 Backbone 的局部策略质量。这一问题指向的不是通信机制本身，而是 Backbone-通信联合训练的稳定性——当 Backbone 本身处于训练动力学的不稳定期时，引入通信支路会放大优化方向的不确定性。训练时较低的测试集规模（16 局）未能及时发现该退化，说明更大规模的评估是必要的。

\textbf{门控自调节的有效性。} gate\_open 在三个种子上均未显著提升胜率，且在 seed 3 上反而低于 normal（$65.6\%$ vs $71.1\%$）。这说明训练学到的门控值（约 $0.06$）并非学习失败——它确实是当前策略状态下的合理融合水平。强制开大门控只会引入额外噪声而不带来收益。

\subsection{通信到动作杠杆的边界诊断}
\esubsection{Boundary Diagnosis of Communication-to-Action Leverage}

在完成 silence budget 与意图对齐实验后，本文进一步进行了两类边界诊断，用于判断后续算法改进究竟应继续围绕“路由/对齐采样”展开，还是应转向“通信如何改变动作 logits”。

\textbf{软冲突对齐的负结果。} 在 hard conflict-aware alignment 中，只有少量不确定或冲突样本被保留，导致梯度高度稀疏且种子敏感。为避免该问题，本文进一步测试了 soft conflict-aware normalized alignment：保留基础有效样本，并使用
\begin{equation}
w_i^t = \mathbb{I}_{\text{valid}}\left(0.2+0.8\max(u_i^t,c_i^t)\right)
\end{equation}
作为软权重，其中 $u_i^t$ 和 $c_i^t$ 分别表示本地不确定性与局部-通信冲突程度，同时损失分母使用基础有效样本数而非被选样本数。该设计成功修复了 hard mask 的样本稀疏问题，但通信隔离评估显示 normal $-$ no\_attack 的平均差值为 $-1.0\%$，并未产生稳定的执行期正向通信贡献。这说明 alignment sampling 的平滑化并不能自动转化为真实动作层收益。

\textbf{门控下限的负结果。} 为检验“通信是否只是因为门控太低而影响不到动作”，本文在 silence budget 主线上加入极轻量的 attack gate floor：将攻击门控下限设为 $0.05$，并用 $0.20$ 作为上界，避免退化为强制开门。结果如表 \ref{tab:leverage_floor} 所示。

\begin{table}[htbp]
    \centering
    \caption{通信强度增强诊断：silence budget 与 gate floor 对比（3 seed 平均）}
    \label{tab:leverage_floor}
    \begin{tabular}{|l|c|c|c|c|c|}
        \hline
        配置 & 最终胜率 & 末 5 点均值 & 峰值 & 攻击门控 & 有效扰动 \\
        \hline
        silence budget & 0.667 & 0.629 & 0.833 & 0.071 & $\approx 6.4{\times}10^{-4}$ \\
        \hline
        gate floor & 0.479 & 0.558 & 0.813 & 0.122 & $1.2{\times}10^{-3}$ \\
        \hline
    \end{tabular}
\end{table}

实验表明，gate floor 确实将攻击门控从 $0.071$ 提升至 $0.122$，有效 logit 扰动也约翻倍；但最终胜率从 $0.667$ 降至 $0.479$，三个种子中有两个种子出现显著退化。与此同时，攻击 no-comm 概率仍保持在约 $0.3\%$--$0.4\%$，攻击熵也维持在健康区间。这说明失败并非来自路由坍缩或沉默逃逸，而是来自残差方向本身尚未被校准：更大的通信注入强度会更稳定地放大未校准扰动，从而伤害已有 Backbone。

\textbf{Advantage-guided leverage 的边界结果。} 基于上述发现，本文进一步测试 advantage-guided attack leverage：仅在正优势攻击样本上，鼓励通信融合后的策略比本地策略更支持实际采取的攻击动作。初始 squared penalty 版本几乎没有激活该机制，attack log-prob 改变量仅约 $1.2\times10^{-5}$，有效扰动约 $4.35\times10^{-4}$。随后改用 Huber penalty 后，机制指标明显增强：攻击门控提升至 $0.171$，有效扰动提升至 $4.95\times10^{-3}$，attack log-prob 改变量提升至 $5.00\times10^{-4}$，且 no-comm 仍被压制在约 $1.5\%$。这说明 Huber leverage 确实让通信 residual 更稳定地进入动作 logits。

\begin{table}[htbp]
    \centering
    \caption{Advantage-guided leverage 诊断结果（3 seed 平均）}
    \label{tab:adv_leverage_huber}
    \begin{tabular}{|l|c|c|c|c|c|c|}
        \hline
        配置 & final & last10 & 攻击门控 & 有效扰动 & $\Delta\log p$ & 隔离增益 \\
        \hline
        silence budget & 0.667 & 0.581 & 0.071 & -- & -- & -- \\
        \hline
        squared leverage & 0.604 & 0.588 & 0.067 & $4.35{\times}10^{-4}$ & $1.18{\times}10^{-5}$ & -- \\
        \hline
        Huber leverage & 0.688 & 0.573 & 0.171 & $4.95{\times}10^{-3}$ & $5.00{\times}10^{-4}$ & $-0.005$ \\
        \hline
    \end{tabular}
\end{table}

然而，对 Huber leverage 的 128 局通信隔离评估给出了更严格的负结果：三 seed 平均 normal 胜率为 $0.4375$，no\_attack 为 $0.4427$，normal $-$ no\_attack 的平均差值为 $-0.52\%$；gate\_open 平均为 $0.4505$，仅比 normal 高 $1.30\%$，且主要来自单个 seed 的小幅波动。同时，gate\_closed 与 no\_attack 在三个 seed 上完全一致，说明诊断开关链路是可信的。由此可见，Huber leverage 证明了“通信可以触达动作 logits”，但没有证明“通信触达的是有因果价值的动作决策”。

\subsection{从动作层杠杆到状态条件化曝光}
\esubsection{From Action-Level Leverage to State-Conditioned Exposure}

上述负结果推动本文进一步重新审视 attack communication 的核心瓶颈。至此已经可以排除几类早期解释：其一，问题不再是 attack stream 完全沉默；其二，问题也不再是 attack attention 完全死亡；其三，问题甚至不再只是门控过低，因为 Huber leverage 已经证明通信 residual 可以稳定进入动作 logits。真正剩下的问题变成了：通信是否在 \emph{该介入的状态} 中改变动作，以及这种改变是否沿着正确方向发生。

基于这一判断，本文后续主线从“继续加大全局通信作用力”转向“只在 peer-local conflict 且 local 较不确定的状态中增强 exposure”。这一主线可以概括为
\begin{equation}
\text{signsplit}
\rightarrow
\text{v4 mild fused exposure}
\rightarrow
\text{v5 margin-near exposure}
\rightarrow
\text{v5b stable exposure}.
\end{equation}

\textbf{Signsplit baseline：通信存在，但几乎不改变动作。} 以 \texttt{v6\_attack\_top2\_selective\_silencebudget\_signsplit\_correction} 为代表的 signsplit 配置已经具备健康的非沉默通信结构：三 seed 平均 no-comm 概率约为 $0.0019$，conflict rate 约为 $0.1130$，说明 peer intent 既存在，也并非总与 local target 重合。然而 peer-vs-local diagnostic 进一步显示，fused follow-peer on conflict rate 为 $0.0000$，fused stay-local on conflict rate 为 $1.0000$，effective delta 仅约 $6\times10^{-4}$。这说明 signsplit 阶段的核心问题不是通信没有被路由出来，而是 communication-to-action leverage 近乎为零。

\textbf{V4 mild fused exposure：第一次把通信稳定推进到动作层。} 为修复 signsplit 的零杠杆问题，本文在 v4 中保留 silence budget，并加入轻量的 peer-conflict margin leverage 与 mild fused exposure。其 peer-local 诊断均值显示：gate mean 从 signsplit 的 $0.0728$ 提升至 $0.5792$，effective delta 从 $0.0006$ 跃升至 $0.0922$，attack-only follow-peer on conflict rate 提升至 $0.0352$，fused follow-peer on conflict rate 也首次达到 $0.0171$。这说明通信已经不再停留在“有路由但碰不到动作”的阶段，而是第一次以可测量的方式进入了 attack action circuit。然而，v4 的 fused stay-local on conflict rate 仍高达 $0.9805$，说明最终决策仍以 local policy 为主导。

\textbf{V5 margin-near exposure：让 fused exposure 更偏向 local 不确定状态。} 在 v4 的基础上，本文进一步引入 local-confidence-aware 的 fused weighting。其核心思想不是盲目让通信更强，而是让 fused 路径在 local margin 更小、local top-1 置信度更低的状态中承受更强的 peer-conflict 压力。相较 v4，这一步使通信介入从“全局 mild exposure”转向“margin-near exposure”，并在部分 seed 上提升了 fused follow-peer 的频率。然而，v5 也暴露出新的不稳定性：最小权重过低时，一些 seed 会重新掉回低 exposure 状态，说明状态条件化虽然是对的方向，但 exposure floor 仍然过弱。

\textbf{V5b stable exposure：将最小不确定性权重从 $0.2$ 提升到 $0.4$。} 为回答“v5 的不稳定是否只是因为 fused exposure floor 太低”这一更窄的问题，本文进一步构造 v5b：保留 margin-near / uncertainty-aware fused exposure 框架，仅将最小不确定性权重由 $0.2$ 提升到 $0.4$，同时保持
\begin{equation}
\texttt{attack\_peer\_conflict\_fused\_local\_conf\_max}=0.90.
\end{equation}
这一步不引入新机制，只测试一个非常小的稳定化修正：在 local 高置信状态中仍保留一部分最小 fused exposure，避免某些 seed 因 fused weight 过低而重新退化到“通信碰不到动作”的 regime。

\subsection{V5b 主线结果与边界}
\esubsection{V5b Mainline Results and Their Boundary}

v5b 的三 seed 通信统计如表 \ref{tab:v5b_three_seed} 所示。

\begin{table}[htbp]
    \centering
    \caption{v5b 三个 seed 的末段通信统计}
    \label{tab:v5b_three_seed}
    \begin{tabular}{|l|c|c|c|c|c|c|c|}
        \hline
        Seed & final win & last5 win & gate last5 & no-comm last5 & entropy last5 & effective delta last5 & fused follow peer last5 \\
        \hline
        1 & 0.6250 & 0.5625 & 0.5364 & 0.0106 & 0.4453 & 0.0930 & 0.0068 \\
        \hline
        2 & 0.9375 & 0.7250 & 0.4416 & 0.0069 & 0.4476 & 0.0740 & 0.0243 \\
        \hline
        3 & 0.6250 & 0.6375 & 0.5642 & 0.0075 & 0.4356 & 0.1077 & 0.0113 \\
        \hline
        Mean & 0.7292 & 0.6417 & 0.5141 & 0.0083 & 0.4428 & 0.0916 & 0.0141 \\
        \hline
    \end{tabular}
\end{table}

这张表有三个关键含义。第一，v5b 维持了健康的非沉默结构：no-comm last5 仅约 $0.0083$。第二，v5b 把 exposure 拉回了强而稳定的区间：gate last5 约为 $0.51$，effective delta last5 约为 $0.092$，基本回到 v4 的强 exposure regime。第三，v5b 并没有把系统推向极端 peer-dominant 决策：fused follow peer last5 只有 $0.0141$，说明最终策略仍然保守。

与 v5 相比，v5b 更像是一个“通信调稳”版本，而不是一个“让通信更激进地主导决策”的版本。其三 seed 平均 final win 从 $0.6042$ 提升至 $0.7292$，last5 win 从 $0.5750$ 提升至 $0.6417$，同时 gate、effective delta 与 fused weight 都明显上升。换言之，v5b 的主要贡献不是让通信更频繁地推翻 local decision，而是让 communication-to-action exposure 不再在不同 seed 上忽强忽弱。

在这一主线上，\texttt{v5b seed2} 是最适合作为论文单条主线展示的 checkpoint。其 final win 达到 $0.9375$，last5 win 为 $0.7250$，last5 return 为 $17.2901$；更重要的是，它的通信参数处于非常干净的区间：gate last5 为 $0.4416$，no-comm last5 为 $0.0069$，entropy last5 为 $0.4476$，effective delta last5 为 $0.0740$。这说明该 checkpoint 并不是靠沉默逃逸、门控饱和或残差爆炸换来的高分，而是在 moderate-gate、healthy-entropy、non-silent 的 regime 中取得了最强结果。

然而，\texttt{v5b seed2} 同时也明确暴露出本文当前最重要的边界。其 fused follow-peer last5 虽提升至 $0.0243$，attack-only follow-peer last5 为 $0.0400$，但 fused stay-local last5 仍高达 $0.9751$，fused log-prob margin 仍显著为负。这意味着通信已经稳定进入动作层，但绝大多数冲突状态下最终攻击决策仍由 local policy 主导。换句话说，v5b 修复的是通信曝光瓶颈（communication exposure bottleneck），而不是最终的因果性决策改变瓶颈（causal decision-change bottleneck）。

基于这一点，本文选择 \texttt{v5b seed2} 作为 paper mainline 的原因并不是“它证明通信已经解决了动作层因果性”，而是它同时满足三个更适合作为论文主线的条件：其一，通信参数健康；其二，通信机制已经 operational at the action level；其三，剩余边界清晰可诊断。这比选择一个更激进但机制不干净的 checkpoint 更适合作为论文的核心叙事。

\subsection{低通信代价主线与 Repo MAIC 对比}
\esubsection{Low-Cost Mainline Versus the In-Repo MAIC}

在明确 \texttt{v5b seed2} 是当前最适合作为论文主线的 checkpoint 之后，本文进一步关心一个更工程化的问题：这一主线是否依赖于较大的执行期通信预算。基于第三章定义的执行期通信代价
\begin{equation}
C_{\mathrm{exec}} = p_{\mathrm{comm}} \cdot k \cdot d_{\mathrm{exec}},
\end{equation}
本文对当前主线与本仓库内的 MAIC 实现进行了粗粒度对比，结果如表 \ref{tab:comm_cost_compare} 所示。

\begin{table}[htbp]
    \centering
    \caption{SMAC \texttt{5m\_vs\_6m} 上执行期通信代价的粗粒度对比}
    \label{tab:comm_cost_compare}
    \begin{tabular}{|p{3.6cm}|p{3.8cm}|c|c|c|c|c|c|}
        \hline
        方法 & 路由形态 & $p_{\mathrm{comm}}$ & $k$ & $d_{\mathrm{exec}}$ & $C_{\mathrm{exec}}$ & $\rho_{\mathrm{topo}}$ & $\rho_{64}$ \\
        \hline
        本文 \texttt{v5b seed2} 攻击流 & top-2 selective + no-comm & 0.9931 & 2 & 8 & 15.9 & 0.4966 & 0.0621 \\
        \hline
        Repo MAIC & dense all-to-all targeted action bias & 1.0000 & 4 & 12 & 48.0 & 1.0000 & 0.1875 \\
        \hline
        稀疏 MAPPO 参考配置 & top-3 sparse comm & 1.0000 & 3 & 4 & 12.0 & 0.7500 & 0.0469 \\
        \hline
        64 维稠密广播参考线 & dense all-to-all hidden-state broadcast & 1.0000 & 4 & 64 & 256.0 & 1.0000 & 1.0000 \\
        \hline
    \end{tabular}
\end{table}

表 \ref{tab:comm_cost_compare} 中，本文主线只统计攻击通信流，因为当前最关键的机制修复、隔离评估与因果诊断都集中发生在该分支上。对 \texttt{v5b seed2} 而言，末段日志给出 attack no-comm probability 约为 $0.0069$，因此
\begin{equation}
p_{\mathrm{comm}} \approx 1 - 0.0069 = 0.9931.
\end{equation}
结合其 top-$2$ 选择性路由和 $8$ 维攻击消息，可得
\begin{equation}
C_{\mathrm{exec}}^{\text{ours, attack}}
\approx
0.9931 \times 2 \times 8
\approx 15.9.
\end{equation}
这意味着在给定单边消息维度的情况下，当前主线只激活了约一半潜在通信边；若进一步以“向其余 $4$ 个队友广播 $64$ 维 hidden state”的稠密方案为参照，则其归一化代价仅为
\begin{equation}
\rho_{64}^{\text{ours, attack}}
=
\frac{15.9}{4\times 64}
\approx 6.2\%.
\end{equation}

相比之下，仓库内 MAIC 的低维 latent 并不能直接视为执行期通信带宽。其配置中的 `latent\_dim=8` 主要服务于内部 teammate modeling，而在执行期真正送入 peer 决策支路的是一个对每个 sender-receiver 对应的动作偏置向量。结合仓库实现可知，MAIC 最终返回的是局部 $Q$ 值与 dense attention 加权 action-bias 的和，因此在 \texttt{5m\_vs\_6m} 上，更合适的执行期消息维度口径是 $d_{\mathrm{exec}} = n_{\mathrm{actions}} = 12$。同时，该实现不包含 top-$k$ 稀疏选择，也不包含显式 no-comm token，因此其执行期通信代价可近似写为
\begin{equation}
C_{\mathrm{exec}}^{\text{MAIC}}
=
1 \times 4 \times 12
= 48.
\end{equation}

据此可得一个对本文叙事十分重要的结论：在相同 \texttt{5m\_vs\_6m} 环境上，当前 \texttt{v5b seed2} 攻击流的执行期通信预算大约仅为 repo MAIC 的三分之一：
\begin{equation}
\frac{15.9}{48} \approx 33.1\%.
\end{equation}
因此，本文 strongest mainline 的优势并不是依赖于高带宽 all-to-all hidden-state communication，而是在显著更低的执行期通信代价下，通过选择性路由、沉默预算和动作层残差注入，获得了稳定的通信机制与当前最强的单条主线结果。

当然，这一对比也应当谨慎理解。首先，表 \ref{tab:comm_cost_compare} 比较的是执行期通信载体，而不是训练期总复杂度，因此没有把 MAIC 的 mutual information 目标和额外推断网络开销计入同一张表。其次，本文当前填写的是攻击流主线预算，而不是完整双流总预算。换言之，这张表支持的不是“本文整体系统绝对最省”，而是一个更克制但更可靠的结论：当前最强主线是在 genuinely low-cost communication regime 中取得的，且其低成本不是表面上的“latent 很小”，而是来自 top-$k$ 稀疏激活、低维消息和受控非沉默概率共同作用的结果。

\subsection{从 V5b 到后续因果修正的边界}
\esubsection{Boundary Beyond V5b}

在 v5b 之后，本文继续尝试了更直接的 bad-action correction 与 causal decision-change 修正，例如在保持 v5b 稳定 exposure 设定的基础上引入更强的 \texttt{bad-action decision change} 分支。其结果表明，仅继续增加 exposure 压力并不会自动把通信变成更好的因果决策器：某些 seed 的 gate 与 effective delta 虽然继续升高，但 final win、last5 win 以及 fused follow-peer 反而下降。这进一步佐证了当前主线的核心判断：更强的通信作用力本身并不是答案，状态选择才是答案。

因此，围绕最新主线可以得出一个比“通信有效/无效”更细的结论：
\begin{enumerate}
    \item 本文已经修复了 attack communication 的沉默逃逸与零曝光问题；
    \item 本文已经让通信 residual 以稳定、可测量的方式进入 attack action logits；
    \item 但本文尚未让通信在大多数冲突状态中稳定跨越 local decision margin；
    \item 因而当前未解的关键问题，不是 communication exposure，而是 causal decision change。
\end{enumerate}

综合来看，silence budget、peer-conflict margin leverage 与 uncertainty-aware fused exposure 共同构成了本文当前最稳的攻击通信主线；而 soft alignment、gate floor、Huber leverage 以及更激进的 bad-action 修正则共同划出了该主线的能力边界。后续算法改进应继续围绕状态条件化的通信因果性建模展开：识别本地策略不确定、peer-local target 冲突、或本地即将执行错误攻击的状态，并让通信只在这些具有潜在因果价值的状态中介入，而不是全局性地让通信“更强”。

\section{本章小结}
\esection{Chapter Summary}

本章首先在 \texttt{join1} 环境中验证了通信机制既可能带来显著收益，也可能由于结构设计不当而完全失效。随后，本文在 SMAC \texttt{5m\_vs\_6m} 环境中构建了较强的 vanilla MAPPO Backbone，系统分析了 Backbone 训练的崩溃-恢复动力学和种子敏感性，并在此基础上发现了冷启动通信的结构性失败。

针对冷启动通信失败，本文进行了 8 组递进式实验，涵盖 IGG、固定门控、通信 warm-up、门控退火和 hidden-state fusion 等多种策略。实验揭示了门控-学习率同构性这一关键洞察，说明门控值在功能上等价于通信学习率。

进一步，本文对攻击通信流进行了系统性机制修复：从 top1 死亡到 top2 软路由、选择性熵正则、沉默预算，最终通过意图对齐将三 seed 平均峰值提升至 $0.746$。在此基础上，本文进一步完成了从 signsplit 到 v4、v5、v5b 的动作层主线推进：先证明通信不再沉默，再证明通信第一次稳定进入动作电路，最后通过 margin-near / uncertainty-aware fused exposure 把跨 seed 的 communication exposure 调稳。其中 \texttt{v5b seed2} 作为当前最适合作为论文主线的 checkpoint，展示了健康的通信参数、稳定的动作层 exposure，以及“causal decision-change 尚未彻底解决”的清晰边界。

这些结果共同表明，通信性能提升并不等价于执行期通信因果贡献的稳定增强。特别是 Huber leverage 与后续更激进的 causal 修正都证明：即使通信已经触达动作 logits，若缺少更精细的状态选择，它仍未必会在正确的时刻沿正确方向改变最终决策。因此，后续改进应继续转向状态条件化的通信因果性建模。

这些发现共同确立了本文的核心论点：轻量、语义明确的通信机制能够产生真实协同价值，但该价值的释放依赖于 Backbone 稳定性、动作层融合方向以及介入状态选择的共同满足。当 Backbone 本身处于训练动力学的不稳定期，或通信残差在不具因果价值的状态中被全局放大时，通信支路会放大优化噪声。这一 insight 将通信研究的焦点从"设计更强的通信"转向"设计更可控、更因果化的通信"。
